\chapter*{Заключение}
\addcontentsline{toc}{chapter}{Заключение}
\label{ch:conclusion}

В ходе выполнения лабораторной работы~\cfgLabNumber{} был разработан
многостраничный макет сайта чемпионата по фиджитал-спорту
\textbf{PHYGITAL~FC} (киберфутбол, Новосибирск-2024) в графическом
редакторе \textbf{Figma}. Спроектированы пять смысловых страниц ---
главная, команда, участник, турнирная таблица и магазин с двумя
секциями (билеты и сувениры), --- каждая из которых выполнена в трёх
адаптивных версиях: десктоп (\(1440\,\)px), планшет (\(768\,\)px) и
мобильный телефон (\(360\,\)px).

Освоенные инструменты и приёмы:
\begin{itemize}
  \item проектирование многостраничного интерфейса с единым
        фирменным стилем (палитра, типографика, иконография);
  \item ускоренный импорт HTML-каркаса в Figma через плагин
        \texttt{html.to.design} с генерацией Auto-Layout-фреймов
        и переменных стилей;
  \item стиль \emph{glassmorphism}: полупрозрачные поверхности с
        размытием, тонкие световые обводки, акцентные неоновые
        цвета на тёмном градиентном фоне;
  \item работа с табличными данными (турнирная таблица,
        статистика игроков) и применение моноширинного шрифта
        \texttt{JetBrains~Mono} для выравнивания цифр;
  \item построение интерактивных компонентов с табами
        (страница магазина), радио-карточек выбора (тип билета)
        и закреплённых блоков (sticky-summary);
  \item настройка \texttt{Constraints} (привязок) для адаптивного
        поведения элементов на разных ширинах экрана;
  \item построение колоночной \texttt{Layout grid} как
        направляющей сетки и подгонка её параметров под три
        размера экрана (\(12/6/4\) колонок);
  \item подготовка адаптивных версий копированием фрейма и
        изменением его ширины с последующей ручной доводкой
        ключевых блоков.
\end{itemize}

Полученные навыки развивают тематику предыдущей лабораторной
работы~10, где макет ограничивался одностраничным лендингом. Опыт
проектирования согласованной системы из пяти страниц с переиспользуемыми
компонентами и единой адаптивной механикой --- основа повседневной
работы веб-дизайнера. Концепции фрейма, привязок и сетки имеют прямые
аналоги в HTML/CSS (\texttt{<div>}, Flexbox/Grid, медиа-запросы), что
облегчает последующий переход от макета к реальной вёрстке сайта
чемпионата.

\endinput
